% !TeX program = xelatex
% 编码：UTF-8。本文仅含第七章，可在当前目录使用 XeLaTeX 编译。
% 排版沿用 MathModel.tex 的 gmcmthesis 类；不含封面、摘要及其他章节。
% 图形直接引用 ../问题4_求解结果/figures/ 中的原始 PDF，无需复制或重画。
% 合并主论文时，取“第七章正文开始”至“第七章正文结束”之间的内容，
% 并在主论文导言区保留下列 graphicspath；表号按主论文原有序列自动衔接。
\documentclass[bwprint]{gmcmthesis}
\usepackage{amsmath}
\usepackage{booktabs}
\usepackage{geometry}
\usepackage{graphicx}
\geometry{a4paper,left=2cm,right=2cm,margin=2.5cm}
\numberwithin{figure}{section}
\renewcommand{\thefigure}{\arabic{section}-\arabic{figure}}
\graphicspath{{../问题4_求解结果/figures/}}
\begin{document}
\setcounter{section}{6}
\setcounter{table}{8}
% 第七章正文开始
\section{问题四的分析与求解}

\subsection{问题描述和分析}

问题四要求在问题三恒定功率协同启动模型的基础上，进一步允许各片电热丝功率随时间独立调节，利用单片温度、电压及其变化信息设计动态辅助加热策略，并在满足冷启动成功与安全约束的条件下，以辅助加热总能耗最小为主要目标，比较不同预冷工况下动态控制与恒功率控制的启动性能。同时，还需建立电堆由常温向低温环境冷却的传热模型，研究预冷时间对初始温度场及后续恒功率冷启动结果的影响。

与前一问相比，本问的关键变化在于初始状态由均匀低温扩展为预冷形成的非均匀温度场，辅助加热则由固定功率输入转变为根据运行状态实时调整的反馈控制。预冷程度决定电堆的初始储热量及端部与中部的温度差异；启动过程中，各片温升、反应产水与结冰状态又随电流加载和辅助功率不断变化。因此，固定功率难以同时适应不同初温及各片的瞬时热量需求，而动态控制需要在增强低温单片加热、抑制冰堵风险与接近目标时及时降功率之间进行协调。

基于上述特点，本文首先建立复合层状电堆的预冷传热模型，将其终态温度场映射为五片单电池及两块端板的初始温度；随后在已有水—冰—热—电耦合模型基础上，引入温度、电压滤波与独立状态观测，构建前馈与PI反馈相结合的分片功率控制律，并通过离线参数搜索确定节能且具有一定安全裕度的控制方案。最后，在统一电流加载、成功判据和评价窗口下比较三种预冷工况的启动结果，并通过预冷时间扫描分析初始储热、片间温差与冷启动性能之间的关系。

\subsection{电堆预冷与动态辅助冷启动模型的建立}

问题四继续采用第五章的七节点电堆热网络和第六章的辅助热源模型，各片内部水冰相变、气体输运、极化损失及电压计算关系保持一致。本问新增的模型包括预冷初始温度场、仅依赖可测信号的状态识别，以及随时间更新的分片辅助功率。为区分两个先后发生的过程，以 $\tau$ 表示预冷时间，以 $t$ 表示电流加载后的启动时间；预冷结束时重新令 $t=0$，并将预冷终态作为冷启动初态。以下物理方程中的温度统一采用开尔文，表格和图中的温度则用摄氏度表示；控制器中的温度差和温升速率不受温标平移影响，使用摄氏度实现时只需将冰点设为零。

\subsubsection{预冷传热模型及初边界条件}

电堆在预冷阶段不加载电流，各片电热丝均关闭，内部不存在电化学反应产热。沿用吹扫后的初始状态，忽略该阶段的水相变及辐射换热，仅考虑沿堆叠方向的导热和端板外表面与环境之间的对流。五片电堆由两块端板、六块双极板及五片MEA组成，其中每片MEA仍分为阳极气体扩散层、阳极催化层、质子交换膜、阴极催化层和阴极气体扩散层，共形成33个实际材料区。内部双极板由相邻单电池共享，不重复计算其厚度和热质量。因此，整堆厚度为
\begin{equation}
L_s=2L_{\mathrm{EP}}+6L_{\mathrm{BP}}+5L_{\mathrm{MEA}}
=2\times10+6\times2+5\times0.3267=33.6335\ \mathrm{mm}.
\tag{7-1}\label{eq:q4:stack_length}
\end{equation}

在第 $j$ 个材料区 $\Omega_j$ 内，采用附件给定的密度、比热容与导热系数，温度场满足
\begin{equation}
(\rho c_p)_j\frac{\partial T}{\partial\tau}
=\frac{\partial}{\partial x}\left(k_j\frac{\partial T}{\partial x}\right),
\qquad x\in\Omega_j,\quad 0<x<L_s.
\tag{7-2}\label{eq:q4:precool_pde}
\end{equation}

相邻材料间取理想热接触，界面处温度和导热通量连续，即
\begin{equation}
T(x_j^-,\tau)=T(x_j^+,\tau),\qquad
k_j\left.\frac{\partial T}{\partial x}\right|_{x_j^-}
=k_{j+1}\left.\frac{\partial T}{\partial x}\right|_{x_j^+}.
\tag{7-3}\label{eq:q4:precool_interface}
\end{equation}

两端对流边界按照向外散热为正确定符号。由于左端外法向与 $x$ 轴正方向相反，左右边界分别为
\begin{equation}
\left\{\begin{aligned}
k_{\mathrm{EP}}\left.\frac{\partial T}{\partial x}\right|_{x=0}
&=h\big[T(0,\tau)-T_{\mathrm{amb}}\big],\\
-k_{\mathrm{EP}}\left.\frac{\partial T}{\partial x}\right|_{x=L_s}
&=h\big[T(L_s,\tau)-T_{\mathrm{amb}}\big],
\end{aligned}\right.
\qquad h=40\ \mathrm{W\,m^{-2}K^{-1}}.
\tag{7-4}\label{eq:q4:precool_robin}
\end{equation}

预冷开始时全堆为均匀常温，环境温度恒为 $-30^\circ\mathrm C$，故初始条件为
\begin{equation}
T(x,0)=T_{\mathrm{init}}=298.15\ \mathrm K,\qquad
T_{\mathrm{amb}}=243.15\ \mathrm K.
\tag{7-5}\label{eq:q4:precool_initial}
\end{equation}

数值计算保留MEA的五层材料结构。为分析电堆热惯性，仍可将其单位面积热容与串联热阻分别合并为
\begin{equation}
\begin{aligned}
C_{A,\mathrm{MEA}}&=\sum_{\ell\in\mathrm{MEA}}(\rho c_p)_\ell L_\ell
=60.75966\ \mathrm{J\,m^{-2}K^{-1}},\\
k_{\mathrm{MEA}}^{\mathrm{eq}}&=\frac{L_{\mathrm{MEA}}}{\displaystyle\sum_{\ell\in\mathrm{MEA}}L_\ell/k_\ell}
\approx0.296\ \mathrm{W\,m^{-1}K^{-1}},\\
C_{A,\mathrm{tot}}&=2C_{A,\mathrm{EP}}+6C_{A,\mathrm{BP}}+5C_{A,\mathrm{MEA}}
=97503.96\ \mathrm{J\,m^{-2}K^{-1}}.
\end{aligned}
\tag{7-6}\label{eq:q4:precool_capacity}
\end{equation}

其中，单块端板与双极板的面热容分别为 $39500$ 和 $3033.36\ \mathrm{J\,m^{-2}K^{-1}}$，两块端板合计约占全堆热容的 $81\%$。因此，预冷温度不仅取决于MEA本身，还受到板体储热及端部散热过程的显著影响，不能仅以单片初始温度的统一平移代替完整预冷计算。

\subsubsection{非均匀初始温度场与电堆热网络的衔接}

预冷模型得到的是沿堆叠方向连续分布的温度场，而启动模型采用五个单电池温度和两个独立端板温度。为保持两种空间描述之间的热量对应关系，将每块内部共享双极板沿几何中面等分，并分别归入左右相邻单电池；两块终端双极板全部归入相应端部单电池，端板则保留独立节点。记第 $a$ 个热节点所属空间范围为 $\Omega_a$，预冷持续时间为 $\tau_c$，则节点初始温度取热容加权平均：
\begin{equation}
\begin{aligned}
T_a^0&=\frac{\displaystyle\int_{\Omega_a}\rho(x)c_p(x)T(x,\tau_c)\,\mathrm dx}
{\displaystyle\int_{\Omega_a}\rho(x)c_p(x)\,\mathrm dx}
=\frac{\displaystyle\sum_{i\in\Omega_a}C_iT_i(\tau_c)}
{\displaystyle\sum_{i\in\Omega_a}C_i},\\
C_i&=(\rho c_p)_i\Delta x_i.
\end{aligned}
\tag{7-7}\label{eq:q4:temperature_map}
\end{equation}

由于各控制体只归属于一个节点，式\eqref{eq:q4:temperature_map}满足预冷物性下的显热恒等关系
\begin{equation}
\sum_{a=1}^{7}C_{A,a}^{\mathrm{pre}}(T_a^0-T_{\mathrm{amb}})
=\sum_iC_i\big[T_i(\tau_c)-T_{\mathrm{amb}}\big],\qquad
C_{A,a}^{\mathrm{pre}}=\sum_{i\in\Omega_a}C_i.
\tag{7-8}\label{eq:q4:map_energy}
\end{equation}

由此构造完全冷却、预冷20 min和预冷40 min三种初态。为便于数值组装，将七节点状态排列为 $\boldsymbol T=[T_1,T_2,T_3,T_4,T_5,T_{\mathrm{EL}},T_{\mathrm{ER}}]^{\mathrm T}$，记上述映射为 $\mathcal M$，则
\begin{equation}
\boldsymbol T^0=\begin{cases}
T_{\mathrm{amb}}\boldsymbol 1, & \text{完全冷却工况},\\
\mathcal M[T(x,1200\ \mathrm s)], & \text{预冷20 min工况},\\
\mathcal M[T(x,2400\ \mathrm s)], & \text{预冷40 min工况}.
\end{cases}
\tag{7-9}\label{eq:q4:three_initial_states}
\end{equation}

启动时各片孔隙液水和冰初值仍取零，膜初始含水量取 $\lambda_0=3$。各片气体、水冰及膜水状态组成内部状态向量 $\boldsymbol X$，随后按第四至第六章的方程演化。辅助功率改为随时间变化后，七节点能量方程可统一写为
\begin{equation}
\begin{aligned}
\boldsymbol C(\boldsymbol X)\frac{\mathrm d\boldsymbol T}{\mathrm dt}
&=-\boldsymbol G(\boldsymbol X)\boldsymbol T
+\boldsymbol hT_{\mathrm{amb}}+\boldsymbol b_{\mathrm{gen}}
+\boldsymbol b_{\mathrm{pc}}+10^4\boldsymbol B\boldsymbol q(t),\\
b_{\mathrm{gen},k}&=10^4j(t)\big(E_{\mathrm{th}}-V_k\big),\qquad
E_{\mathrm{th}}=1.48\ \mathrm V,\qquad k=1,\ldots,5.
\end{aligned}
\tag{7-10}\label{eq:q4:plant_heat}
\end{equation}

其中，$\boldsymbol B$ 将五片电热丝功率输入对应单电池节点，其端板行均为零；$\boldsymbol G$ 包含片间、片板导热及环境散热，$\boldsymbol b_{\mathrm{pc}}$ 表示水相变化的热效应。$j$ 与 $q_k$ 分别以 $\mathrm{A\,cm^{-2}}$ 和 $\mathrm{W\,cm^{-2}}$ 计，因而进入单位面积热网络时均须乘以 $10^4$。

这里保留了前章启动模型的两项继承设置：启动阶段采用含水气有效物性的节点热容，对流散热作用于首尾单电池节点；预冷阶段则采用附件材料热容，对流作用于实际端板外表面。因此，式\eqref{eq:q4:map_energy}保证预冷温度场在其自身物性下的热量加权一致性，但不意味着两阶段具有完全相同的热容或严格相同的初始焓。这一处理使不同辅助控制策略能够在同一启动模型上比较，其近似影响通过后续传热参数扰动与数值检验加以考察。

\subsubsection{温度、电压反馈与单电池运行状态识别}

控制器可直接获取的反馈量为各片温度和电压，冰含量及端板温度则需通过模型估计。为避免将仿真中的真实内部状态作为控制输入，另设具有独立水、气和热状态的名义模型观测器。观测器由已知电流与上一时段实际施加的加热功率推进，每个采样时刻仅用测得的五片温度校正对应温度节点；其端板温度及内部水气状态保留模型预测值。记控制时刻 $t_n=n\Delta t_c$，可写为
\begin{equation}
\begin{aligned}
\widehat{\boldsymbol x}^{\,n,-}
&=\mathcal P_{\widehat{\boldsymbol p}}\!\left(
\widehat{\boldsymbol x}^{\,n-1,+},j|_{[t_{n-1},t_n]},\boldsymbol q^{n-1},\Delta t_c\right),
\\
\widehat T_k^{n,+}&=T_{k,\mathrm m}^n,\qquad k=1,\ldots,5.
\end{aligned}
\tag{7-11}\label{eq:q4:independent_observer}
\end{equation}

式中，$\widehat{\boldsymbol p}$ 为名义物性参数，$\mathcal P$ 表示水—热—电耦合模型的预测算子；实际电堆中的冰、水、气库存和扰动物性均不传入观测器。其预测电压 $\widehat V_k^n$ 由独立内部状态、采样温度与已知电流共同计算，作为后续冰风险修正的依据。

为减弱采样噪声对控制动作的影响，首先对测温和测压信号进行一阶低通滤波：
\begin{equation}
\begin{aligned}
\widetilde T_k^n&=\beta\widetilde T_k^{n-1}+(1-\beta)T_{k,\mathrm m}^n,\\
\widetilde V_k^n&=\beta\widetilde V_k^{n-1}+(1-\beta)V_{k,\mathrm m}^n,
\qquad \beta=\exp\!\left(-\frac{\Delta t_c}{0.15\ \mathrm s}\right).
\end{aligned}
\tag{7-12}\label{eq:q4:signal_filter}
\end{equation}

再由滤波信号的后向差分估计温升速率与电压变化率，并对差分结果继续滤波：
\begin{equation}
\begin{aligned}
r_{T,k}^n&=\beta_r r_{T,k}^{n-1}+(1-\beta_r)
\frac{\widetilde T_k^n-\widetilde T_k^{n-1}}{\Delta t_c},\\
r_{V,k}^n&=\beta_r r_{V,k}^{n-1}+(1-\beta_r)
\frac{\widetilde V_k^n-\widetilde V_k^{n-1}}{\Delta t_c},\qquad
\beta_r=\exp\!\left(-\frac{\Delta t_c}{0.20\ \mathrm s}\right).
\end{aligned}
\tag{7-13}\label{eq:q4:rate_filter}
\end{equation}

设 $\widehat\varepsilon_{k,\mathrm{mod}}^n$ 为独立模型给出的片内最大冰体积分数，采用预测电压与测量电压的有符号残差修正冰风险标量：
\begin{equation}
\begin{aligned}
\widehat\varepsilon_k^n
&=\operatorname{sat}_{[0,1]}\!\left[
\widehat\varepsilon_{k,\mathrm{mod}}^n
+L_\varepsilon\frac{\widehat V_k^n-V_{k,\mathrm m}^n}{\delta V_{\mathrm{obs}}}\right],
\\
&\qquad L_\varepsilon=0.02,\quad\delta V_{\mathrm{obs}}=0.10\ \mathrm V.
\end{aligned}
\tag{7-14}\label{eq:q4:ice_estimator}
\end{equation}

其中，$\operatorname{sat}_{[a,b]}(z)=\min\{b,\max(a,z)\}$。当实测电压低于模型预测时，修正项提高冰风险估计；反向残差则降低估计值。该修正只作用于控制器使用的风险标量，不直接改写观测器的水冰库存，也不等价于从电压唯一重建真实冰场。

以 $T_f=273.15\ \mathrm K$ 表示冰点，$T_{\mathrm{targ}}$ 和 $t_{\mathrm{plan}}$ 分别表示参考目标温度和规划升温时间。温度偏低程度、当前所需温升速率及温升不足程度定义为
\begin{equation}
\begin{aligned}
D_{T,k}^n&=\max\!\left(0,\frac{T_f-\widetilde T_k^n}{30\ \mathrm K}\right),\\
r_{\mathrm{req},k}^n&=\max\!\left(0,\frac{T_{\mathrm{targ}}-\widetilde T_k^n}
{\max(t_{\mathrm{plan}}-t_n,1\ \mathrm s)}\right),\\
D_{r,k}^n&=\max\!\left(0,\frac{r_{\mathrm{req},k}^n-r_{T,k}^n}
{r_{\mathrm{req},k}^n+0.01\ \mathrm{K\,s^{-1}}}\right).
\end{aligned}
\tag{7-15}\label{eq:q4:temperature_deficit}
\end{equation}

进一步将冰量估计、电压裕度和电压下降速率组合为风险指标
\begin{equation}
\begin{aligned}
R_k^n={}&0.4\frac{\widehat\varepsilon_k^n}{\varepsilon_{\mathrm{crit}}}
+0.4\max\!\left(0,\frac{V_{\mathrm{safe}}+\delta_V-\widetilde V_k^n}{\delta_V}\right)
\\
&{}+0.2\min\!\left(1,\frac{\max(0,-r_{V,k}^n)}{0.10\ \mathrm{V\,s^{-1}}}\right).
\end{aligned}
\tag{7-16}\label{eq:q4:ice_risk}
\end{equation}

其中，$V_{\mathrm{safe}}=0.30\ \mathrm V$、$\varepsilon_{\mathrm{crit}}=0.99$ 为安全阈值，$\delta_V=0.10\ \mathrm V$ 为电压预警带宽。当冰量、电压、下降速率或综合风险任一项触发预警时，判为需要优先增强加热：
\begin{equation}
\begin{aligned}
\mathcal D_k^n={}&
\big\{\widehat\varepsilon_k^n>0.8\big\}
\ \cup\ \big\{\widetilde V_k^n<0.40\ \mathrm V\big\}
\ \cup\ \big\{R_k^n>0.7\big\}\\
&{}\cup\ \big\{r_{V,k}^n<-0.10\ \mathrm{V\,s^{-1}},\
\widetilde V_k^n<0.60\ \mathrm V\big\}.
\end{aligned}
\tag{7-17}\label{eq:q4:danger_condition}
\end{equation}

上述判据将低温但温升正常、温升不足、冰堵或低压风险以及接近目标等情形区分开来，使控制器能够依据当前状态分配辅助功率。其作用是提前响应风险，是否真正满足题设安全约束仍须由被控对象的真实状态独立检验。

\subsubsection{前馈与PI结合的动态功率控制律}

动态控制以参考温升轨迹为基本目标，在初期弥补升温所需热量，在临近目标时降低输入，并在低压或冰堵预警时保留安全增强通道。各片采用自身初温建立分段线性参考轨迹：
\begin{equation}
\begin{aligned}
T_{\mathrm{ref},k}(t)&=T_k^0+
(T_{\mathrm{targ}}-T_k^0)\min\!\left(\frac{t}{t_{\mathrm{plan}}},1\right),
\\
\dot T_{\mathrm{ref},k}(t)&=\begin{cases}
(T_{\mathrm{targ}}-T_k^0)/t_{\mathrm{plan}}, & t<t_{\mathrm{plan}},\\
0, & t\geq t_{\mathrm{plan}}.
\end{cases}
\end{aligned}
\tag{7-18}\label{eq:q4:temperature_reference}
\end{equation}

目标温度和规划时间作为离线搜索参数，不预先假定某一轨迹具有全局最优性。以滤波温度构造跟踪误差，并根据低温程度调整比例增益：
\begin{equation}
e_k^n=T_{\mathrm{ref},k}(t_n)-\widetilde T_k^n,
\qquad
K_{P,k}^n=K_P\big[0.6+0.4\min(1,D_{T,k}^n)\big].
\tag{7-19}\label{eq:q4:tracking_error_gain}
\end{equation}

前馈项利用独立模型的节点热容、导热损失、估计端板温度及反应产热，补偿沿参考轨迹升温时的净热量缺口：
\begin{equation}
\begin{aligned}
q_{\mathrm{ff},k}^n
&=\frac{\alpha_{\mathrm{ff}}}{10^4}
\left[\widehat C_{A,k}^n\dot T_{\mathrm{ref},k}(t_n)
+\widehat\ell_k^n-10^4j(t_n)\big(E_{\mathrm{th}}-V_{k,\mathrm m}^n\big)\right],
\\
\widehat{\boldsymbol\ell}^{\,n}
&=\widehat{\boldsymbol G}^{\,n}\widehat{\boldsymbol T}^{\,n,+}
-\widehat{\boldsymbol h}T_{\mathrm{amb}}.
\end{aligned}
\tag{7-20}\label{eq:q4:feedforward}
\end{equation}

其中，$\widehat\ell_k^n$ 表示对应电池节点通过热网络净流出的面热流，$\alpha_{\mathrm{ff}}$ 为前馈系数，常规分支取1。式\eqref{eq:q4:feedforward}显式扣除了反应产热，使电流加载带来的升温贡献能够减少辅助热需求；相变及模型误差造成的偏差由反馈项进一步修正。

控制器采用位置式PI，并对积分量实施限幅和饱和方向判断。先计算试探积分量及相应未限幅输出：
\begin{equation}
\begin{aligned}
I_{k,\mathrm{tr}}^n&=I_k^{n-1}+K_Ie_k^n\Delta t_c,\\
u_{k,\mathrm{tr}}^n&=q_{\mathrm{ff},k}^n+K_{P,k}^ne_k^n+I_{k,\mathrm{tr}}^n
+K_r\min(D_{r,k}^n,2)\min(D_{T,k}^n,1).
\end{aligned}
\tag{7-21}\label{eq:q4:trial_integral}
\end{equation}

令 $q_{\max}=1\ \mathrm{W\,cm^{-2}}$、$I_{\max}=2\ \mathrm{W\,cm^{-2}}$。仅当积分更新不会继续加剧上下限饱和时接受该增量，即
\begin{equation}
I_k^n=\begin{cases}
\operatorname{sat}_{[-I_{\max},I_{\max}]}(I_{k,\mathrm{tr}}^n),
& \begin{aligned}
&(u_{k,\mathrm{tr}}^n\leq q_{\max}\ \text{或}\ e_k^n<0)\\[-2pt]
&\quad\text{且}\ (u_{k,\mathrm{tr}}^n\geq0\ \text{或}\ e_k^n>0),
\end{aligned}\\
I_k^{n-1}, & \text{其他情况}.
\end{cases}
\tag{7-22}\label{eq:q4:antiwindup}
\end{equation}

由此得到名义功率
\begin{equation}
q_{\mathrm{nom},k}^n=q_{\mathrm{ff},k}^n+K_{P,k}^ne_k^n+I_k^n
+K_r\min(D_{r,k}^n,2)\min(D_{T,k}^n,1).
\tag{7-23}\label{eq:q4:nominal_power}
\end{equation}

最后一项为温升不足补偿，$K_r$ 为相应增益。该项同时受低温程度和温升不足程度调节，因而主要在冰点以下且升温滞后时增加加热，避免仅依据温度误差而忽略升温速度。前馈每次直接参与当前名义功率计算，不向上一时刻功率重复累加。

当滤波温度已越过冰点、估计冰量较低且电压不再快速下降时，按温度距目标的剩余量设置渐缩因子：
\begin{equation}
\psi_k^n=\begin{cases}
\displaystyle\operatorname{sat}_{[0,1]}\!\left(
\frac{T_{\mathrm{targ}}-\widetilde T_k^n}{\Delta T_{\mathrm{taper}}}\right),
&\begin{aligned}
&\widetilde T_k^n>T_f,\ \widehat\varepsilon_k^n<0.5,\\[-2pt]
&r_{V,k}^n>-0.05\ \mathrm{V\,s^{-1}},
\end{aligned}\\
1, &\text{其他情况}.
\end{cases}
\tag{7-24}\label{eq:q4:taper_factor}
\end{equation}

其中，$\Delta T_{\mathrm{taper}}$ 为渐缩温度带宽。对于式\eqref{eq:q4:danger_condition}识别的风险状态，则由综合风险计算增强功率下限：
\begin{equation}
\begin{aligned}
q_{\mathrm{boost},k}^n={}&q_{\mathrm b}
+(q_{\max}-q_{\mathrm b})\operatorname{sat}_{[0,1]}\!\left(
\frac{R_k^n-R_{\mathrm w}}{1-R_{\mathrm w}}\right),
\\
&\qquad q_{\mathrm b}=0.6\ \mathrm{W\,cm^{-2}},\quad R_{\mathrm w}=0.7.
\end{aligned}
\tag{7-25}\label{eq:q4:safety_boost}
\end{equation}

最终执行功率先对名义输入限幅并渐缩，再与安全增强下限取大值，保证近目标降功率不会抵消已经触发的风险响应：
\begin{equation}
\begin{aligned}
q_k^n&=\operatorname{sat}_{[0,q_{\max}]}\!\left[
\max\!\left\{\psi_k^n\operatorname{sat}_{[0,q_{\max}]}(q_{\mathrm{nom},k}^n),
\ \boldsymbol 1_{\mathcal D_k^n}q_{\mathrm{boost},k}^n\right\}\right],
\\
q_k(t)&=q_k^n,\qquad t_n\leq t<t_{n+1}.
\end{aligned}
\tag{7-26}\label{eq:q4:applied_power}
\end{equation}

式中，$\boldsymbol 1_{\mathcal D_k^n}$ 为风险状态指示量。各片独立执行相同形式的反馈律，但由于测量温度、电压及估计状态不同，其输出功率可以不同。控制状态与相应动作归纳于表\ref{tab:q4:states}。关热锁存一旦完成即覆盖所有加热动作；在尚未关热时，安全增强具有最高优先级。该优先级设计能够避免逻辑冲突，但有限功率下是否有足够升温能力仍取决于实际工况，不能仅凭控制结构保证任意扰动下均满足安全约束。

\begin{table}[htbp]
    \centering
    \caption{单电池运行状态识别及动态辅助加热动作}
    \vspace{0.3cm}
    \label{tab:q4:states}
    \renewcommand{\arraystretch}{1.4}
    \small
    \setlength{\tabcolsep}{4pt}
    \begin{tabular}{c p{2.7cm} p{6.0cm} p{4.8cm}}
        \toprule
        \textbf{状态码} & \textbf{运行状态} & \textbf{识别依据} & \textbf{控制动作} \\
        \midrule
        1 & 低压或冰堵风险 & 冰量、电压、下降速率或综合风险触发式\eqref{eq:q4:danger_condition} & 在渐缩后施加安全功率下限，优先增强加热 \\
        2 & 温升不足 & 未触发更高优先级动作，且 $D_{r,k}>0.1$ & 前馈与PI跟踪，低温时叠加温升不足补偿 \\
        3 & 正常跟踪 & 未触发风险及渐缩条件，且温升不足程度较低 & 按参考轨迹调节名义功率 \\
        4 & 接近目标 & 满足式\eqref{eq:q4:taper_factor}中的低冰量、稳定电压及越过冰点条件 & 逐步降低名义功率，风险触发时仍执行增强 \\
        5 & 测量保持完成 & 全部单片连续满足测量停机条件2 s & 五片功率统一置零并永久锁存 \\
        \bottomrule
    \end{tabular}
\end{table}

\subsubsection{辅助加热能耗、安全约束与停机判据}

设实际关热时刻为 $t_{\mathrm{off}}$，第 $k$ 片辅助能耗以及电堆总辅助能耗分别为
\begin{equation}
E_{\mathrm{aux},k}=A\int_0^{t_{\mathrm{off}}}q_k(t)\,\mathrm dt,
\qquad E_{\mathrm{aux}}=\sum_{k=1}^{5}E_{\mathrm{aux},k},
\qquad A=25\ \mathrm{cm^2}.
\tag{7-27}\label{eq:q4:auxiliary_energy}
\end{equation}

由于 $q_k$ 以 $\mathrm{W\,cm^{-2}}$ 计，式\eqref{eq:q4:auxiliary_energy}使用面积的 $\mathrm{cm^2}$ 数值后直接得到焦耳。能耗只计电热丝消耗的电能，电化学反应热另行记账。加热功率、电压及冰量约束要求在整个主评价窗口内均成立：
\begin{equation}
\left\{\begin{aligned}
&0\leq q_k(t)\leq1\ \mathrm{W\,cm^{-2}},\\
&\min_{k,\,0\leq t\leq t_{\mathrm{off}}}V_k(t)\geq0.30\ \mathrm V,\\
&\max_{k,\,x,\,0\leq t\leq t_{\mathrm{off}}}\varepsilon_{\mathrm{ice},k}(x,t)<0.99.
\end{aligned}\right.
\tag{7-28}\label{eq:q4:path_constraints}
\end{equation}

其中，冰体积分数沿用前章以局部总体积为基准的定义，不与孔隙冰饱和度混用。除题设安全量外，还检查局部加载电流小于极限电流、有效气相孔隙率非负、质子电导为正及各相库存非负，以排除虽满足温度条件但内部物理状态已经失效的数值解。

为与问题三形成可比结果，五片单电池采用相同的给定电流加载，累计电荷可直接积分为
\begin{equation}
\begin{aligned}
j(t)&=\min(0.005t,0.3)\ \mathrm{A\,cm^{-2}},\\
Q(t)&=\int_0^tj(s)\,\mathrm ds
=0.0025[\min(t,60)]^2+0.3\max(t-60,0)
\quad\big(\mathrm{C\,cm^{-2}}\big),
\end{aligned}
\tag{7-29}\label{eq:q4:load_charge}
\end{equation}

式中时间数值以秒计。本文延续问题三的启动电荷比较预算 $Q(t_s)\leq20\ \mathrm{C\,cm^{-2}}$，由前60 s消耗电荷 $9\ \mathrm{C\,cm^{-2}}$ 可得首次成功期限
\begin{equation}
t_s\leq t_Q=60+\frac{20-9}{0.3}=96.6667\ \mathrm s.
\tag{7-30}\label{eq:q4:startup_deadline}
\end{equation}

该预算用于统一前后两问的比较边界，不另行解释为问题四新增的题设约束。为区分物理启动结果与控制器的测量判断，记 $\mathcal P(t)$ 表示截至时刻 $t$ 的真实电压、冰量及内部物理有效性均满足要求，则真实首次成功时刻定义为
\begin{equation}
t_s=\inf\left\{t\geq0:\ \min_{1\leq k\leq5}T_k(t)>T_f,
\quad\mathcal P(t)\ \text{成立}\right\}.
\tag{7-31}\label{eq:q4:true_first_success}
\end{equation}

式\eqref{eq:q4:true_first_success}用于评价和事件定位，不直接触发关热。控制器仅在采样时刻检查可获得的滤波温度、滤波电压与冰量估计，并留出 $0.2\ \mathrm K$ 的测温裕度：
\begin{equation}
\begin{aligned}
b_n=\boldsymbol 1\big\{&
\min_k\widetilde T_k^n>T_f+0.2\ \mathrm K,
\quad\min_k\widetilde V_k^n\geq0.30\ \mathrm V,
\\
&\max_k\widehat\varepsilon_k^n<0.99\big\}.
\end{aligned}
\tag{7-32}\label{eq:q4:sensor_condition}
\end{equation}

设 $H_n$ 为连续满足测量条件的已确认时长，初值 $H_0=0$。当任一次采样不满足条件时清零；连续两个采样均满足时才累计采样间隔：
\begin{equation}
H_n=\begin{cases}
H_{n-1}+\Delta t_c, & b_n=b_{n-1}=1,\\
0, & \text{其他情况},
\end{cases}
\qquad n\geq1.
\tag{7-33}\label{eq:q4:hold_counter}
\end{equation}

当全部单片连续满足测量条件达到 $t_{\mathrm{hold}}=2\ \mathrm s$ 后，统一关闭五片电热丝并永久锁存，即
\begin{equation}
t_{\mathrm{off}}=\min\{t_n:H_n\geq t_{\mathrm{hold}}\},
\qquad q_k(t)=0,\quad t\geq t_{\mathrm{off}},\quad k=1,\ldots,5.
\tag{7-34}\label{eq:q4:shutdown_latch}
\end{equation}

在每次仿真结束后，另以真实物理状态验证关热前是否确实连续成功至少2 s；若测量条件提前成立而真实保持不满足要求，则该组参数判为不可行。因此，$t_s$ 表示真实状态首次达到启动要求的时刻，$t_{\mathrm{off}}$ 表示测量保持确认后的实际关热时刻，二者通常不满足简单的 $t_{\mathrm{off}}=t_s+2\ \mathrm s$。首次事件须满足式\eqref{eq:q4:startup_deadline}，常规测量保持窗口最多延至 $t_Q+t_{\mathrm{hold}}$；$Q(t_s)$ 与 $Q(t_{\mathrm{off}})$ 分别保存，避免把保持阶段的额外加载混入首次启动预算。若预冷初态已温暖且初始真实状态与可用测量均通过检查，则直接记 $t_s=t_{\mathrm{off}}=0$，不施加冷启动辅助加热。

用于策略主比较的最低电压、最大冰体积分数和电堆最大温差，均在 $[0,t_{\mathrm{off}}]$ 内统计：
\begin{equation}
\begin{aligned}
V_{\min}&=\min_{0\leq t\leq t_{\mathrm{off}}}\min_kV_k(t),\\
\varepsilon_{\max}&=\max_{0\leq t\leq t_{\mathrm{off}}}\max_{k,x}\varepsilon_{\mathrm{ice},k}(x,t),\\
\Delta T_{\max}&=\max_{0\leq t\leq t_{\mathrm{off}}}
\left[\max_kT_k(t)-\min_kT_k(t)\right].
\end{aligned}
\tag{7-35}\label{eq:q4:comparison_metrics}
\end{equation}

其中，温差取同一时刻五片单电池的最高与最低温度之差，再对时间取最大值，不将不同时刻的温度极值相减，也不混入端板温度。真实首次成功窗口 $[0,t_s]$ 的能耗及极值另行记录，用于与第六章原有口径对照。关热后继续按给定电流加载观察60 s，其温度、电压和冰量作为持续运行特性的后验结果单独分析；完成启动及2 s保持并不自动意味着关热后始终维持在冰点以上。

\subsection{动态加热控制策略优化与数值求解}
在 7.2 建立预冷模型与动态辅助加热控制律后，需要进一步确定参考温升轨迹及反馈增益，使五片电堆在满足启动与安全约束的条件下，以尽可能小的辅助电能完成冷启动。由于功率限幅、状态切换、测量保持和温度过零事件均会改变控制响应，能耗关于控制参数一般不光滑，难以直接采用单一起点的梯度方法求解。因此，本文沿用前两章的“内层瞬态仿真—外层参数优化”框架，在离线仿真中整定控制参数，在线阶段仅根据温度、电压反馈计算分片加热功率，并进一步通过扰动场景筛选获得具有启动裕度的推荐方案。

\subsubsection{控制参数、优化目标与可行域}
动态控制律虽然对五片电池分别计算功率，但各片采用相同的控制结构与参数，功率差异由各片温度、电压、温升速率及估计热负荷决定。因此，无需逐时刻独立优化五条功率曲线，而将常规搜索变量写为
\begin{equation}
\boldsymbol{\theta}=\left(T_{\mathrm{targ}},t_{\mathrm{plan}},K_P,K_I,\Delta T_{\mathrm{taper}},K_r\right).
\tag{7-36}\label{eq:q4:solver_parameters}
\end{equation}
其中，\(K_r\) 为温升不足补偿增益，其余参数分别控制目标温度、参考升温时间、比例反馈、积分反馈及近目标降功率区间。各参数的搜索范围见表\ref{B7-control-bounds}。\(t_{\mathrm{plan}}\) 是参考轨迹的设计参数，可以超过实际启动期限；候选方案是否可行仍由真实首次成功时刻与实际关热时刻判定，不能以参考轨迹的终点代替启动成功。

\begin{table}[htbp]
    \centering
    \caption{动态辅助加热控制参数的常规搜索范围}
    \vspace{0.3cm}
    \label{B7-control-bounds}
    \renewcommand{\arraystretch}{1.5}
    \begin{tabular}{l c c c}
        \toprule
        控制参数 & 单位 & 下界 & 上界 \\
        \midrule
        目标温度 \(T_{\mathrm{targ}}\) & \(^\circ\mathrm C\) & 0.4 & 5 \\
        规划时间 \(t_{\mathrm{plan}}\) & \(\mathrm s\) & 15 & 150 \\
        比例增益 \(K_P\) & \(\mathrm{W\,cm^{-2}\,K^{-1}}\) & 0.015 & 0.5 \\
        积分增益 \(K_I\) & \(\mathrm{W\,cm^{-2}\,K^{-1}\,s^{-1}}\) & 0.001 & 0.08 \\
        渐缩带宽 \(\Delta T_{\mathrm{taper}}\) & \(\mathrm K\) & 0.1 & 2 \\
        温升不足补偿增益 \(K_r\) & \(\mathrm{W\,cm^{-2}}\) & 0 & 0.4 \\
        \bottomrule
    \end{tabular}
\end{table}

为使策略比较具有一致性，采样周期、测量停机规则及安全阈值不参与上述六维搜索。名义搜索中前馈倍率固定为 1，连续保持时间取 \(2\,\mathrm s\)，安全增强起始功率取 \(0.6\,\mathrm{W\,cm^{-2}}\)，冰量预警阈值、电压预警带和综合风险阈值分别取 0.8、\(0.1\,\mathrm V\) 和 0.7，电压创新校正增益取 0.02。另设关闭前馈、比例、积分和温升不足补偿的独立候选，用于检验当前初温下能否仅依靠反应热启动；该候选保留安全后备动作，不受表\ref{B7-control-bounds}中反馈增益正下界的限制。

记满足 7.2 所述功率、真实状态、测量保持与物理有效性条件的参数集合为 \(\Theta_{\mathrm{feas}}\)。首次真实成功需不晚于 \(96.6667\,\mathrm s\)，测量确认关热需在 \(98.6667\,\mathrm s\) 的有限评价窗口内完成。辅助能耗统计至实际关热时刻，归一化目标为
\begin{equation}
\min_{\boldsymbol{\theta}\in\Theta_{\mathrm{feas}}}
J(\boldsymbol{\theta})
=\frac{E_{\mathrm{aux}}(\boldsymbol{\theta})}{E_{\mathrm{ref}}}.
\tag{7-37}\label{eq:q4:solver_objective}
\end{equation}
式中，\(E_{\mathrm{ref}}\) 取同工况固定恒功率方案的辅助能耗，在同一次搜索中保持不变，仅用于量纲归一化。电压、冰量和内部状态有效性直接作为可行性条件处理；搜索时对不可行候选附加高罚值，但最终选择不允许以能耗收益抵偿安全约束违背。

为避免极小的能耗差掩盖明显的启动时间差，正式精度复算后，在有限可行候选集合 \(\mathcal C_f\) 中先确定最低辅助能耗，再构造其 \(0.05\%\) 近优带：
\begin{equation}
E_{\min}=\min_{\boldsymbol{\theta}\in\mathcal C_f}E_{\mathrm{aux}}(\boldsymbol{\theta}),\qquad
\mathcal C_{\varepsilon}=\left\{\boldsymbol{\theta}\in\mathcal C_f:
E_{\mathrm{aux}}(\boldsymbol{\theta})\leq1.0005E_{\min}\right\}.
\tag{7-38}\label{eq:q4:solver_energy_band}
\end{equation}
在该集合中，依次优先选择首次真实成功时间较短、最大同刻片间温差较小和辅助能耗较低的候选。用于排序的温差采用式\eqref{eq:q4:comparison_metrics}中的 \(\Delta T_{\max}\)，比较同一时刻五片单电池的温度，能够表征启动过程中的片间一致性；若将不同时刻的最高温度与最低温度相减，则会把整体升温幅度混入空间温差，不能用于本问的一致性评价。

\subsubsection{预冷离散与闭环耦合求解方法}
预冷阶段按照实际层状结构划分 33 个材料区，共设置 146 个控制体，材料界面与控制体边界对齐。设第 \(i\) 个控制体宽度、导热系数及单位面积热容分别为 \(\Delta x_i\)、\(k_i\) 和 \(C_{A,i}\)，则相邻控制体之间及边界控制体至环境之间的单位面积热导为
\begin{equation}
g_{i+1/2}=\left(\frac{\Delta x_i}{2k_i}+\frac{\Delta x_{i+1}}{2k_{i+1}}\right)^{-1},\qquad
g_b=\left(\frac{\Delta x_b}{2k_b}+\frac{1}{h}\right)^{-1}.
\tag{7-39}\label{eq:q4:solver_precool_conductance}
\end{equation}
边界热导中同时保留半个控制体的导热阻力与外表面对流阻力，使离散散热方向与预冷模型的 Robin 边界一致。令 \(\boldsymbol{u}=\boldsymbol{T}-T_{\mathrm{amb}}\boldsymbol{1}\)，以 \(\mathbf C_{\mathrm p}\) 表示对角热容矩阵、\(\mathbf L_{\mathrm p}\) 表示由内部与边界热导组装的矩阵，则后向欧拉格式为
\begin{equation}
\left(\frac{\mathbf C_{\mathrm p}}{\Delta\tau}+\mathbf L_{\mathrm p}\right)
\boldsymbol{u}^{m+1}=\frac{\mathbf C_{\mathrm p}}{\Delta\tau}\boldsymbol{u}^{m}.
\tag{7-40}\label{eq:q4:solver_precool_be}
\end{equation}
正式预冷时间步取 \(\Delta\tau=0.25\,\mathrm s\)，并在所需预冷时刻精确落点。得到空间温度场后，按照 7.2 的热容加权映射生成五片单电池和两个端板的初始温度，保留端板与单电池之间的初始热差异。

启动阶段每片采用 58 个基础传输控制体描述水、冰及气体状态，热状态仍由七节点网络计算。为区别于控制采样时刻 \(t_n\)，将内部积分网格记为 \(s_m\)，相邻步长为 \(\Delta s_m=s_{m+1}-s_m\)。给定当前加热功率后，先推进质量迁移与相变，再通过热—电 Picard 迭代求解温度、电压和反应产热，其热网络离散形式为
\begin{equation}
\begin{aligned}
\left(\frac{\mathbf C^{m+1}}{\Delta s_m}+\mathbf G^{m+1}\right)
\boldsymbol{T}^{m+1}
={}&\frac{\mathbf C^{m+1}}{\Delta s_m}\boldsymbol{T}^{m}
+\boldsymbol b_{\mathrm{amb}}^{m+1}+\boldsymbol q_{\mathrm{gen}}^{m+1}\\
&+\frac{\Delta\boldsymbol H_{\mathrm{phase}}^{m+1}}{\Delta s_m}
+\boldsymbol q_{\mathrm{aux}}(s_m).
\end{aligned}
\tag{7-41}\label{eq:q4:solver_startup_be}
\end{equation}
其中，各热源均按单位面积计量，\(\Delta\boldsymbol H_{\mathrm{phase}}\) 为当前步的相变释热量，辅助源向量的前五个分量为 \(10^4q_k\)，两个端板分量为零。正式积分步长取 \(0.025\,\mathrm s\)，Picard 迭代的温度增量收敛容差取 \(10^{-9}\,\mathrm K\)，每步最多迭代 30 次，并记录实际迭代误差用于最终可行性复核。

控制器每隔 \(0.2\,\mathrm s\) 采集一次温度和电压，经滤波及独立状态观测后更新加热功率；相邻采样点之间功率保持不变。观测器用名义物性、已知电流及已施加的功率独立推进，其内部水冰状态不从被控对象复制。为使数值积分与控制更新一致，积分步在采样时刻、电流平台切换时刻和评价窗口终点截短。步内电流采用累计电荷的精确增量确定，即
\begin{equation}
\bar j_m=\frac{Q(s_{m+1})-Q(s_m)}{\Delta s_m},\qquad
E_{\mathrm{aux}}=25\sum_{s_m<t_{\mathrm{off}}}\Delta s_m\sum_{k=1}^{5}q_k(s_m).
\tag{7-42}\label{eq:q4:solver_charge_energy}
\end{equation}
式中，\(q_k(s_m)\) 表示当前控制采样区间保持的执行功率，其单位为 \(\mathrm{W\,cm^{-2}}\)，因此面积取 \(25\,\mathrm{cm^2}\) 后能耗单位为焦耳。

当一个积分步跨越五片共同过零事件时，从该步起点重新积分，并以二分法定位首次真实成功时刻；该事件仅用于评价，不提前触发测量停机。停机仍由采样时刻的滤波温度、电压及估计冰量连续满足规定条件 \(2\,\mathrm s\) 后决定，随后永久锁存零功率。每个内部时间步同时检查电压、冰量、剩余气相孔隙率、物种库存、质子电导率及极限电流条件，防止仅凭温度越零或测量保持完成将物理无效轨迹判为成功。

\subsubsection{离线参数搜索与名义节能方案求解}
对于完全冷却、预冷 20 min 和预冷 40 min 三种初态，分别调用相同的闭环仿真器进行参数整定。首先评估关闭名义加热的独立候选。若其实际辅助功率始终为零且满足全部可行性条件，则辅助能耗达到非负下界，无需继续以降低名义辅助能耗为目标搜索。预冷 20 min 工况属于这一情形，但零辅助能耗仅说明该名义初态能够自启动，不能据此推断其具有足够的扰动裕度。

对于需要辅助加热的工况，先对已有物理初值和默认参数进行多起点复核，再以两个固定随机种子运行差分进化搜索。初筛阶段采用 \(0.1\,\mathrm s\) 时间步和每片 29 个传输控制体，以降低大量候选评估的计算成本；随后提取能耗较低且参数互异的 16 个可行候选，统一以 \(0.025\,\mathrm s\) 时间步、每片 58 个控制体重新计算，并以复算后的较优候选为起点进行有界 Powell 局部整定。最终仅在正式精度下的可行候选中，按照式\eqref{eq:q4:solver_energy_band}规定的能耗近优带和排序规则选取名义方案。

上述流程同时覆盖快速升温与延长加热时间两类可能的解，并通过正式精度重排排除粗步长下出现的虚假可行候选。由于实际搜索限于给定反馈结构、参数范围及有限候选集合，所得方案应理解为当前搜索范围内的优选可行解，不能将某次差分进化结果解释为任意连续功率函数空间中的严格全局最优解。

\subsubsection{扰动场景筛选与留裕度控制方案确定}
名义节能方案可能靠近首次成功期限，在初温降低或端部散热增强时失去可行性。因此，在名义参数附近进一步构造留裕度候选：对完全冷却和预冷 40 min 工况，提高目标温度并适当缩短规划时间；对预冷 20 min 工况，恢复小幅比例、积分及温升不足补偿，并比较不同前馈倍率，以获得必要的主动加热储备。

每个候选均在相同的六个训练场景中计算，场景联合考虑初温 \(\pm1\,\mathrm K\) 偏移、片间热导、端板耦合热导及对流系数的 \(\pm20\%\) 偏差，并叠加标准差为 \(0.2\,\mathrm K\) 的温度噪声和 \(0.005\,\mathrm V\) 的电压噪声。除共同物理与停机条件外，训练筛选还要求
\begin{equation}
t_s(\boldsymbol\theta,\omega)\leq94.6667\,\mathrm s,\qquad
t_{\mathrm{off}}(\boldsymbol\theta,\omega)\leq94.6667\,\mathrm s,
\quad \forall\omega\in\Omega_{\mathrm{train}}.
\tag{7-43}\label{eq:q4:solver_training_deadline}
\end{equation}
该要求相对于首次成功期限保留 \(2\,\mathrm s\) 裕度，相对于允许关热的最终窗口保留 \(4\,\mathrm s\) 裕度。训练初筛采用 \(0.05\,\mathrm s\) 时间步和每片 58 个控制体；从通过全部训练场景的候选中，按名义能耗选取前四个，再以正式步长 \(0.025\,\mathrm s\) 复核名义工况及六个训练场景，最终取复核通过者中名义辅助能耗最低的方案。以下将名义节能方案记为 \(D_{\mathrm N}\)，推荐留裕度方案记为 \(D_{\mathrm R}\)，三种工况的对应参数见表\ref{B7-control-parameters}。

\begin{table}[htbp]
    \centering
    \caption{三种预冷工况的名义节能与推荐留裕度控制参数}
    \vspace{0.3cm}
    \label{B7-control-parameters}
    \renewcommand{\arraystretch}{1.5}
    \resizebox{\textwidth}{!}{
    \begin{tabular}{l l c c c c c c c}
        \toprule
        工况 & 方案 & \(T_{\mathrm{targ}}/{}^\circ\mathrm C\) & \(t_{\mathrm{plan}}/\mathrm s\) & \(K_P\) & \(K_I\) & \(\Delta T_{\mathrm{taper}}/\mathrm K\) & \(K_r\) & 前馈倍率 \\
        \midrule
        完全冷却 & \(D_{\mathrm N}\) & 0.45227 & 97.15711 & 0.49466 & 0.02328 & 0.26215 & 0.36643 & 1 \\
        完全冷却 & \(D_{\mathrm R}\) & 3.50000 & 95.21397 & 0.49466 & 0.02328 & 0.26215 & 0.36643 & 1 \\
        预冷 20 min & \(D_{\mathrm N}\) & 1.00000 & 60.00000 & 0 & 0 & 1.00000 & 0 & 0 \\
        预冷 20 min & \(D_{\mathrm R}\) & 1.00000 & 95.00000 & 0.05000 & 0.00500 & 1.00000 & 0.05000 & 0 \\
        预冷 40 min & \(D_{\mathrm N}\) & 0.73357 & 98.57205 & 0.21677 & 0.03518 & 0.60698 & 0.20018 & 1 \\
        预冷 40 min & \(D_{\mathrm R}\) & 2.00000 & 96.60061 & 0.21677 & 0.03518 & 0.60698 & 0.20018 & 1 \\
        \bottomrule
    \end{tabular}}
\end{table}

由表\ref{B7-control-parameters}可见，完全冷却和预冷 40 min 工况主要通过提高目标温度并略微提前参考轨迹增加裕度，而预冷 20 min 工况由零名义加热转为小增益反馈补偿。推荐方案因此承担额外的辅助能耗，其选择目标是在给定训练场景下保留可行性，不应与纯名义条件下的最低能耗混为一谈。

参数确定后予以冻结，再进行独立验证。每种工况在三种初温偏移与十组新噪声种子的组合下开展 30 次试验，三种工况共 90 次；另对三个传热参数分别施加正负偏差，并设置两个联合偏差场景，每种工况 8 次，共 24 次物性失配与噪声联合试验。训练与独立验证使用不同随机种子，后者不再用于逐次调参。相应通过率及失败机理在 7.4 中分析，其结论限于已列出的扰动范围与样本，不能外推为任意扰动下的严格安全保证。

\subsubsection{恒功率对照与预冷时间扫描设置}
为回答动态策略相对于问题三恒功率协同策略的改进效果，首先保留问题三的固定功率分配
\begin{equation}
\overline{\boldsymbol q}_{C}=
(1,\,1,\,0.6245115588,\,1,\,1)\,\mathrm{W\,cm^{-2}}.
\tag{7-44}\label{eq:q4:solver_constant_power}
\end{equation}
三工况主比较中，将采用测量保持规则的固定恒功率方案记为 \(C_{\mathrm H}\)，其与动态策略使用相同初态、电流曲线、采样滤波和测量保持规则，均以真实首次成功时刻评价启动速度，以实际关热时刻统计辅助电能及启动过程极值。另保留在真实首次成功时立即关热的理想恒功率结果 \(C_0\)，用于复现问题三基准；该结果的统计窗口不同，不与具有测量保持的能耗直接混用。

为区分动态控制收益与重新整定功率分配的收益，进一步对每种初态单独优化恒功率方案，记为 \(C_*\)。先在镜像形式 \((q_e,q_n,q_c,q_n,q_e)\) 下进行差分进化和 Powell 搜索，再加入五片非对称邻域扰动及低功率候选扫描，并以正式精度复算。该对照与动态主方案采用相同保持条件，另外在 \(30\,\mathrm s\)、\(60\,\mathrm s\)、\(90\,\mathrm s\) 及 \(96.6667\,\mathrm s\) 的首次成功期限下构造有限候选能耗—时间前沿。这样可以同时考察固定旧分配、同工况重新优化恒功率与推荐动态控制之间的差异，但恒功率名义再优化结果不代表其经过同等扰动整定后的能力上限。

针对预冷程度的影响，固定采用式\eqref{eq:q4:solver_constant_power}的功率向量，将预冷时间设置为
\begin{equation}
\tau_c\in\{10,15,20,\ldots,100\}\,\mathrm{min}.
\tag{7-45}\label{eq:q4:solver_cooling_scan}
\end{equation}
对 19 个预冷时刻分别求解空间温度场、映射启动初态并进行恒功率冷启动仿真，记录启动时间、辅助能耗、最低电压、最大冰体积分数及最大同刻片间温差。该扫描沿用问题三的理想首次成功口径，统计窗口截止于 \(t_s\)，不附加测量保持时间；若初态已满足成功条件，则直接记录 \(t_s=0\) 和零辅助能耗。由此可在固定加热方案下单独辨识预冷程度的作用，避免将初始温度变化与控制参数重新整定的影响混合。

最后，对冻结参数后的正式方案分别开展预冷空间与时间加密、启动步长与传输网格加密、控制采样周期变化以及水量和热量守恒检查。三工况主方案还在关热后继续施加题给电流曲线观察 \(60\,\mathrm s\)，该段始终保持零辅助功率，其温度、电压和冰量极值单独统计，用于判断首次启动和短时保持完成后能否继续维持暖态。

\subsection{优化结果与动态辅助冷启动特性分析}
\label{sec:q4:results}

在上述预冷模型与闭环求解方法基础上，分别计算完全冷却、预冷20 min和预冷40 min三种工况。为兼顾辅助能耗与已测试扰动下的可行性，本节以留裕度动态方案\(D_R\)作为推荐策略，以沿用问题三功率分配并采用相同测量保持规则的恒功率方案\(C_H\)作为主要对照。两者使用相同的初始温度场、电流加载曲线和测量停机条件，比较结果见表\ref{tab:q4:main}。表中启动时间为真实状态首次满足成功判据的时刻\(t_s\)，辅助电能及电压、冰量和温差极值均统计至实际关热时刻\(t_{\rm off}\)；关热后的延续运行结果另行讨论。

\begin{table}[htbp]
    \centering
    \caption{不同预冷工况下动态与恒功率辅助启动结果对比}
    \vspace{0.3cm}
    \label{tab:q4:main}
    \renewcommand{\arraystretch}{1.5}
    \resizebox{\textwidth}{!}{
    \begin{tabular}{l c c c c c c c c c}
    \toprule
    \textbf{工况} & \textbf{策略} & \shortstack{\textbf{功率控制}\\\textbf{方式}} & \shortstack{\textbf{启动时间}\\\(t_s/\mathrm s\)} & \shortstack{\textbf{关热时刻}\\\(t_{\rm off}/\mathrm s\)} & \shortstack{\textbf{辅助电能}\\\(E_{\rm aux}/\mathrm J\)} & \shortstack{\textbf{最大温差}\\\(\Delta T_{\max}/\mathrm K\)} & \shortstack{\textbf{最低电压}\\\(V_{\min}/\mathrm V\)} & \shortstack{\textbf{最大局部}\\\textbf{冰体积分数}} & \textbf{结果}\\
    \midrule
    完全冷却 & \(C_H\) & 固定分片功率 & 25.741 & 28.2 & 3260.281 & 27.003 & 0.66185 & 0.012130 & 成功\\
    完全冷却 & \(D_R\) & 分片动态反馈 & 85.244 & 88.0 & 2929.996 & 6.060 & 0.56044 & 0.277253 & 成功\\
    预冷20 min & \(C_H\) & 固定分片功率 & 5.606 & 8.0 & 924.902 & 8.716 & 0.73621 & \(9.528\times10^{-5}\) & 成功\\
    预冷20 min & \(D_R\) & 分片动态反馈 & 74.308 & 78.2 & 293.609 & 9.054 & 0.60497 & 0.087293 & 成功\\
    预冷40 min & \(C_H\) & 固定分片功率 & 17.037 & 19.4 & 2242.888 & 19.852 & 0.68486 & 0.002818 & 成功\\
    预冷40 min & \(D_R\) & 分片动态反馈 & 88.539 & 91.6 & 1855.385 & 7.243 & 0.57332 & 0.239568 & 成功\\
    \bottomrule
    \end{tabular}}
\end{table}

由表\ref{tab:q4:main}可知，三种工况下两类策略均完成启动，但其能量投入和时间分配明显不同。动态策略将三工况辅助电能分别降至2929.996 J、293.609 J和1855.385 J，同时将启动过程延长至74--89 s，使反应热在升温过程中承担更大比例的热负荷。因此，动态策略的节能效果需要结合启动速度、片间一致性及结冰水平共同评价，不能仅根据辅助电能降低便认为全部性能同步改善。

\subsubsection{三种预冷工况的初始温度分布特征}

预冷结束时的温度场决定了后续启动所需补充的显热，同时影响各片电池与端板之间的初始传热方向。表\ref{tab:q4:initial}给出三种工况的代表节点温度和不同空间范围的温差。由于左右结构及边界近似对称，将两端对应节点合并列示；表中热容加权均温包含端板、双极板及MEA，全场温差包含端板真实外表面，五片温差则仅由五个单电池平均温度计算。

\begin{table}[htbp]
    \centering
    \caption{三种工况的初始温度及空间非均匀性}
    \vspace{0.3cm}
    \label{tab:q4:initial}
    \renewcommand{\arraystretch}{1.5}
    \resizebox{\textwidth}{!}{
    \begin{tabular}{l c c c c c c c}
    \toprule
    \textbf{工况} & \shortstack{\textbf{端板均温}\\\(/{}^\circ\mathrm C\)} & \shortstack{\textbf{第1、5片}\\\(/{}^\circ\mathrm C\)} & \shortstack{\textbf{第2、4片}\\\(/{}^\circ\mathrm C\)} & \shortstack{\textbf{第3片}\\\(/{}^\circ\mathrm C\)} & \shortstack{\textbf{热容加权均温}\\\(/{}^\circ\mathrm C\)} & \shortstack{\textbf{全场温差}\\\(/\mathrm K\)} & \shortstack{\textbf{五片温差}\\\(/\mathrm K\)}\\
    \midrule
    完全冷却 & -30.0000 & -30.0000 & -30.0000 & -30.0000 & -30.0000 & 0 & 0\\
    预冷20 min & -9.2554 & -9.0879 & -8.9790 & -8.9493 & -9.2123 & 0.50698 & 0.13860\\
    预冷40 min & -22.1591 & -22.0958 & -22.0547 & -22.0434 & -22.1428 & 0.19162 & 0.05239\\
    \bottomrule
    \end{tabular}}
\end{table}

\begin{figure}[!h]
\centering
\includegraphics[width=.95\textwidth]{图01_预冷温度场族.pdf}
\caption{不同预冷时间下的整堆温度场及局部温差}
\label{F7-1}
\end{figure}

由图\ref{F7-1}(a)及表\ref{tab:q4:initial}可见，预冷20 min和40 min时温度均呈中部较高、两端较低的分布，但其均温分别约为\(-9.21^\circ\mathrm C\)和\(-22.14^\circ\mathrm C\)，仍明显高于完全冷却工况。端板既承担对外散热，又占据整堆的大部分热容，因此控制了整体降温的主要时间尺度。离散热系统的最慢模态时间常数约为1233.26 s，即20.55 min，与总热容除以两端换热导度所得的集总估算接近。预冷从20 min延长至40 min后，电堆与环境之间的温差继续减小，整堆也逐步接近均匀的\(-30^\circ\mathrm C\)状态。

图\ref{F7-1}(b)将各时刻温度减去相应最低温度后显示局部梯度，可以看到MEA所在薄层的温度变化较集中，而高导热双极板内部温度变化较小。在附件给定物性和端面换热条件下，20 min的全场温差约为0.507 K，40 min时反而减小至0.192 K；对应的五片平均温差仅为0.139 K和0.052 K。因此，两种未完全冷却工况确实存在空间非均匀性，但不能由预冷时间延长直接推出绝对温差增大。本模型给出的主要变化是整体初温进一步降低，初始梯度在较长预冷阶段逐渐衰减。后续启动中若出现数开尔文乃至数十开尔文的片间温差，其主要来源应从端板吸热、分片加热和反应产热的动态竞争中解释。

\subsubsection{动态功率调节与温度、电压及结冰演化}

图\ref{F7-2}--图\ref{F7-4}分别给出三种工况推荐策略的功率、温度、电压和局部冰峰值。各图竖直分界线表示测量保持完成并锁存关热的时刻，阴影部分为关热后60 s的独立观察窗口。因此，评价表\ref{tab:q4:main}中的主结果时应使用分界线之前的数据，而不能把后验段的冰量增长或温度回落混入主窗口极值。

\begin{figure}[!h]
\centering
\includegraphics[width=.95\textwidth]{图04_case1_动态控制轨迹.pdf}
\caption{完全冷却工况下推荐动态策略的功率与水热电状态演化}
\label{F7-2}
\end{figure}

对于完全冷却工况，图\ref{F7-2}(a)显示，各片在启动初期均获得短时辅助热输入，随后内部第2--4片的功率逐渐降至零，而端部第1、5片的加热功率随热负荷变化持续增加，后期达到\(1\ \mathrm{W\,cm^{-2}}\)上限。内部电池能够利用自身反应热及片间热传递继续升温，端部电池则同时承担端板吸热与环境散热，因而成为决定全堆首次越过冰点的限制位置。动态控制通过减少内部电池的持续加热，将主要辅助电能集中于端部，使过程最大片间温差控制在6.060 K，显著低于固定恒功率方案的27.003 K。

\begin{figure}[!h]
\centering
\includegraphics[width=.95\textwidth]{图05_case2_动态控制轨迹.pdf}
\caption{预冷20 min工况下推荐动态策略的功率与水热电状态演化}
\label{F7-3}
\end{figure}

预冷20 min时，各部件仍保留较多显热，辅助加热需求明显减小。由图\ref{F7-3}(a)可见，端部功率先缓慢增加再逐渐回落，峰值仅为\(0.1243\ \mathrm{W\,cm^{-2}}\)，内部各片功率更低。端部电池累计实际加热时间为67.4 s，此后辅助功率已经降至零，而全堆首次达标时刻为74.308 s，测量保持确认时刻为78.2 s。这表明该工况末段主要依靠反应热和已形成的热状态完成升温，控制器无需在整个启动窗口持续供热。端部第1、5片在初期还存在短暂回落，说明较高初温虽然降低了总体热量缺口，但加载初期反应热较弱时，环境散热和端板吸热仍可能超过即时热输入。

\begin{figure}[!h]
\centering
\includegraphics[width=.95\textwidth]{图06_case3_动态控制轨迹.pdf}
\caption{预冷40 min工况下推荐动态策略的功率与水热电状态演化}
\label{F7-4}
\end{figure}

预冷40 min的初温介于前两种工况之间，图\ref{F7-4}(a)所示端部功率总体低于完全冷却工况，峰值约为\(0.5978\ \mathrm{W\,cm^{-2}}\)；内部第2--4片仅在早期加热约18.4 s，端部加热则维持至91.6 s。虽然该工况辅助电能低于完全冷却工况，但其首次启动时间88.539 s略长于完全冷却工况的85.244 s，原因是两种初态分别采用独立整定的控制参数，具有不同的功率水平和参考升温轨迹。因此，跨工况比较时，初温高低能够解释热量需求的变化，却不能脱离控制参数直接判断启动时间的排序。

为了量化空间功率分配，表\ref{tab:q4:cellenergy}给出推荐策略的逐片辅助电能。完全冷却、预冷20 min和预冷40 min工况中，端部两片合计分别占总辅助电能的约96.5\%、91.5\%和97.0\%。与问题三固定策略对第2、4片长期施加满功率不同，动态方案能够根据实际温升及时降低内部片功率，从而减少内部过热。这里的功率分配是片间热耦合与反馈共同作用的结果，并非预先规定只有端部允许加热。

\begin{table}[htbp]
    \centering
    \caption{推荐动态策略的逐片辅助电能分配}
    \vspace{0.3cm}
    \label{tab:q4:cellenergy}
    \renewcommand{\arraystretch}{1.5}
    \begin{tabular}{l c c c c c c}
    \toprule
    \textbf{工况} & \textbf{第1片/J} & \textbf{第2片/J} & \textbf{第3片/J} & \textbf{第4片/J} & \textbf{第5片/J} & \textbf{合计/J}\\
    \midrule
    完全冷却 & 1413.246 & 34.529 & 34.447 & 34.529 & 1413.246 & 2929.996\\
    预冷20 min & 134.347 & 9.674 & 5.566 & 9.674 & 134.348 & 293.609\\
    预冷40 min & 899.927 & 18.566 & 18.399 & 18.566 & 899.927 & 1855.385\\
    \bottomrule
    \end{tabular}
\end{table}

从图\ref{F7-2}--图\ref{F7-4}的电压和冰量子图可见，电流爬升过程中反应产水与低温冻结共同增加冰量，而温度升高和部分区域越过冰点后又促进融冰。推荐策略三工况的最低单片电压分别为0.56044 V、0.60497 V和0.57332 V，局部最大冰体积分数分别为0.277253、0.087293和0.239568，均满足0.30 V和0.99的约束。相比固定恒功率，动态策略在低温区停留更久、累计加载电荷更大，因此主窗口冰峰值更高，最低电压也更低。由此可见，本次优化是在保持安全可行的条件下减少辅助电能，并未同时实现最强的抑冰效果。

还需区分控制结构中设置的安全保护与实际轨迹中触发的动作。状态记录表明，三种推荐方案在无扰动主工况下均未进入安全增强状态，主要控制过程由温升跟踪、温升不足补偿及近目标降功率构成。安全增强提供的是风险上升时的后备规则，不能将其描述为这些主工况节能或抑冰的已发生原因。图中冰量采用全MEA各控制体的局部冰体积分数最大值，与MEA体积平均冰量及孔隙冰饱和度含义不同，后两者不能替代本节的安全约束指标。

\subsubsection{辅助能耗、启动速度与恒功率策略综合比较}

在相同测量保持规则下，定义推荐动态方案相对固定恒功率的辅助节电率为
\begin{equation}
\eta_E=\frac{E_{\rm aux}^{C_H}-E_{\rm aux}^{D_R}}{E_{\rm aux}^{C_H}}\times100\%.
\tag{7-46}\label{eq:q4:saving}
\end{equation}
由表\ref{tab:q4:main}计算得到，完全冷却、预冷20 min和预冷40 min的辅助节电率分别为10.13\%、68.26\%和17.28\%。其中20 min工况初始热量储备较高，沿用完全冷却工况整定的固定大功率会产生较强的加热冗余，动态控制降低功率后的相对收益最大；完全冷却工况需补充的显热较多，端部功率后期达到上限，因此节电空间相对有限。

上述辅助节电伴随着明显的时间代价。固定恒功率三工况的首次启动时间分别为25.741 s、5.606 s和17.037 s，而推荐动态方案分别为85.244 s、74.308 s和88.539 s。推荐方案首次达标时的累计电荷依次为16.573、13.293和17.562\(\ \mathrm{C\,cm^{-2}}\)，虽然均未超过继承的20\(\ \mathrm{C\,cm^{-2}}\)预算，但明显高于固定恒功率的1.657、0.079和0.726\(\ \mathrm{C\,cm^{-2}}\)。更长的反应时间提供了更多反应热，同时也意味着消耗更多反应物。因此，式\eqref{eq:q4:saving}仅衡量电热丝辅助电能的变化，不能直接解释为氢耗或全系统总能耗的同幅下降。

片间一致性也存在工况差异。完全冷却和40 min工况的最大温差分别由27.003 K和19.852 K降至6.060 K和7.243 K，表明动态控制有效抑制了内部片过热；但20 min工况由8.716 K略增至9.054 K。较暖初态允许控制器长时间采用较小辅助功率，端部热损失与内部反应热的不均衡仍会逐步形成温差。因此，动态策略对温差的改善主要体现在本次两个较冷工况，不能概括为所有初态下一致性均优于恒功率。

\begin{figure}[!h]
\centering
\includegraphics[width=.95\textwidth]{图14_重优化恒功率与推荐动态.pdf}
\caption{固定恒功率、同工况重新优化恒功率与推荐动态策略的比较}
\label{F7-5}
\end{figure}

为进一步区分反馈调节本身的收益与原固定功率分配的适用性，图\ref{F7-5}补充了按各工况重新整定的恒功率方案\(C_*\)。在相同测量停机规则下，\(C_*\)三工况能耗分别为3237.876 J、0 J和1984.470 J，首次启动时间分别为29.416 s、84.725 s和95.287 s。推荐动态方案在完全冷却和40 min工况下相对\(C_*\)分别减少307.880 J和129.085 J，节电约9.51\%和6.50\%；其中40 min工况还将首次启动提前约6.75 s，但完全冷却工况仍需更长时间。

20 min工况的重新优化恒功率退化为五片零加热，名义动态优选方案也得到相同的0 J结果。这说明该初态在给定模型及加载条件下能够利用反应热实现自启动，辅助能耗的非负下界已经达到。推荐方案投入293.609 J，用于缩短启动时间并提高已测试扰动下的可行性，因此不能将其称为该工况的名义最小能耗，也不能将相对原固定功率的全部节电归因于反馈结构独有的优势。重新优化恒功率仅进行了名义参数搜索，其结果也不能代表经过同等扰动整定后的恒功率能力上限。

此外，本章固定功率的首次达标复算能耗约为2976.015 J，而采用测量保持规则后增至3260.281 J。前者与第六章约2975.3 J的结果接近，微小差异与数值设置及事件求解精度有关；后者增加的能耗主要来自温度裕度确认与保持期间继续加热。因此，本章节电率始终使用表\ref{tab:q4:main}中的共同停机基准，避免将第六章首次达标能耗与本章实际关热能耗直接相除。

\subsubsection{预冷时间对恒功率冷启动及单片一致性的影响}

针对问题四第（2）小问，将预冷时间依次取为10、15、\(\ldots\)、100 min，并始终施加问题三的固定功率分配和统一电流曲线。该组计算用于隔离预冷程度的影响，采用首次真实达标即结束的口径，不附加测量温度裕度和2 s保持。图\ref{F7-6}给出全部19个扫描点，表\ref{tab:q4:scan}选取其中六个代表时刻列出结果。

\begin{figure}[!h]
\centering
\includegraphics[width=.95\textwidth]{图08_预冷时间扫描.pdf}
\caption{预冷时间对固定恒功率启动性能的影响}
\label{F7-6}
\end{figure}

\begin{table}[htbp]
    \centering
    \caption{不同预冷时间下固定恒功率启动的代表性结果}
    \vspace{0.3cm}
    \label{tab:q4:scan}
    \renewcommand{\arraystretch}{1.5}
    \resizebox{\textwidth}{!}{
    \begin{tabular}{c c c c c c c c}
    \toprule
    \shortstack{\textbf{预冷时间}\\\(/\mathrm{min}\)} & \shortstack{\textbf{热容加权初温}\\\(/{}^\circ\mathrm C\)} & \shortstack{\textbf{首次达标}\\\(/\mathrm s\)} & \shortstack{\textbf{辅助电能}\\\(/\mathrm J\)} & \shortstack{\textbf{最大五片温差}\\\(/\mathrm K\)} & \shortstack{\textbf{最低电压}\\\(/\mathrm V\)} & \shortstack{\textbf{最大局部}\\\textbf{冰体积分数}} & \textbf{结果}\\
    \midrule
    10 & 3.812 & 0 & 0 & 0.225 & 1.23732 & 0 & 初态已达标\\
    20 & -9.212 & 5.606 & 648.136 & 6.298 & 0.74738 & \(9.528\times10^{-5}\) & 成功\\
    40 & -22.143 & 17.037 & 1969.681 & 17.469 & 0.68684 & 0.002818 & 成功\\
    60 & -27.030 & 22.316 & 2579.978 & 22.416 & 0.67089 & 0.007411 & 成功\\
    80 & -28.878 & 24.428 & 2824.169 & 24.161 & 0.66534 & 0.010141 & 成功\\
    100 & -29.576 & 25.242 & 2918.317 & 24.799 & 0.66329 & 0.011347 & 成功\\
    \bottomrule
    \end{tabular}}
\end{table}

由图\ref{F7-6}(a)--(c)可见，预冷越充分，整堆初温越接近\(-30^\circ\mathrm C\)，固定功率启动所需时间和辅助电能也随之增加，但增加幅度逐渐减小。10 min时五片初温约为4.01--4.24\(^\circ\mathrm C\)，起始状态已经满足温度条件，因此首次达标时间和辅助能耗均为零；该点不代表经历了一个零时间的加热过程，也不构成任意后续负载下可持续运行的证明。20 min和40 min的辅助电能分别为648.136 J和1969.681 J，而100 min时达到2918.317 J，逐步接近完全冷却工况约2976.015 J的首次达标基准。由于所有扫描点的总加热功率保持不变，能耗增长与首次启动时间延长直接对应，其根源是预冷后需要补充的显热增加。

图\ref{F7-1}与图\ref{F7-6}(d)还揭示了初始温差和启动温差的不同变化方向。随着预冷时间由10 min增加到100 min，完整预冷场的绝对温差由约0.825 K减至0.010 K，五片初始均温差由约0.225 K减至0.003 K；但启动过程最大片间温差却由0.225 K增加至24.799 K。较深预冷使大热容端板处于更低温度，启动时端部电池需要持续向端板供热，而内部电池在固定加热和反应产热共同作用下较快升温。因而，较均匀的低温初态并不意味着较均匀的启动过程，启动期间重新形成的端部热负荷差异才是后期片间不一致性的主要原因。

在电压与冰量方面，预冷时间增加使低温产水持续时间延长，局部冰峰值逐渐升高，最低单片电压逐渐降低。即使预冷100 min，最低电压仍为0.66329 V，局部冰峰值仅为0.011347，均保持在规定阈值内。全部19个离散预冷时间点均完成首次启动，说明问题三固定功率方案在所考察区间内具有可行性；本次扫描没有发现失败点，因此不能据此给出未经识别的临界预冷时间。该结论仅针对已计算的19个点及固定参数条件。

\subsubsection{关热后温度回落与持续运行特性}

首次越过冰点并完成短时保持，只说明启动阶段满足规定判据，并不等价于整个电堆已经具备维持暖态的热储备。为检验这一差异，在各方案实际关热后继续按题给电流曲线运行60 s，辅助功率保持锁存为零。表\ref{tab:q4:post}给出推荐方案关热时的状态及后验段极值，同时列出固定恒功率方案的后验最低温度作为参照。

\begin{table}[htbp]
    \centering
    \caption{推荐动态策略关热状态及关热后60 s的运行结果}
    \vspace{0.3cm}
    \label{tab:q4:post}
    \renewcommand{\arraystretch}{1.5}
    \resizebox{\textwidth}{!}{
    \begin{tabular}{l c c c c c c}
    \toprule
    \textbf{工况} & \shortstack{\textbf{关热最低片温}\\\(/{}^\circ\mathrm C\)} & \shortstack{\textbf{关热左端板温度}\\\(/{}^\circ\mathrm C\)} & \shortstack{\textbf{后验最低片温}\\\(/{}^\circ\mathrm C\)} & \shortstack{\textbf{后验最低电压}\\\(/\mathrm V\)} & \shortstack{\textbf{后验局部}\\\textbf{冰峰值}} & \shortstack{\(C_H\)\textbf{后验最低片温}\\\(/{}^\circ\mathrm C\)}\\
    \midrule
    完全冷却 & 0.882 & -14.471 & -4.508 & 0.55967 & 0.50104 & -7.986\\
    预冷20 min & 0.433 & -5.115 & 0.433 & 0.61376 & 0.08702 & -2.677\\
    预冷40 min & 0.763 & -10.288 & -1.910 & 0.58853 & 0.29566 & -6.556\\
    \bottomrule
    \end{tabular}}
\end{table}

由表\ref{tab:q4:post}及图\ref{F7-2}、图\ref{F7-4}的阴影部分可见，完全冷却和40 min工况虽然在关热时全部单电池均已高于冰点，但端板仍分别约为\(-14.47^\circ\mathrm C\)和\(-10.29^\circ\mathrm C\)。辅助热源一旦撤去，端板继续从端部单电池吸热，使最低片温分别降至\(-4.508^\circ\mathrm C\)和\(-1.910^\circ\mathrm C\)。这与第六章发现的端板热惯性效应一致，说明仅根据单电池温度过零及2 s保持，尚不能判定端板热汇已经充分减弱。

相比之下，20 min工况的初始热储备更高，且在正式关热之前辅助功率已提前降至零，反应热能够维持其继续升温；关热后最低片温保持在0.433\(^\circ\mathrm C\)以上。完全冷却和40 min工况虽再次进入零下区间，但后验最低电压仍高于0.30 V，冰峰值也未达到0.99，因此应表述为未维持持续暖态，而不是将其与电压失效或严重冰堵混为一谈。与固定恒功率方案相比，推荐动态方案在本次60 s观察窗口内减轻了回冷，但仍未在全部工况消除该现象。

上述后验段不计入主表辅助电能，且不允许重新开启加热。若进一步要求关热后始终保持冰点以上，则需将这一要求明确纳入优化约束，并重新整定功率与热裕度；本次结果不能直接外推为已满足该更强目标。后验段继续加载也会消耗反应物，其作用是检验延续运行状态，不能混入首次启动的电荷预算。

\subsubsection{参数灵敏度与控制策略稳健性分析}

图\ref{F7-7}分别展示预冷传热参数变化和名义动态控制参数扰动的影响。预冷部分同时考察绝对温差及其相对于整堆平均超温的比值，以区分内部传热不均匀程度与整体冷却速度；启动部分则固定名义控制结构，对片间导度、端板耦合导度、环境换热及控制参数逐项扰动。

\begin{figure}[!h]
\centering
\includegraphics[width=.95\textwidth]{图11_物理参数与控制器敏感性.pdf}
\caption{预冷传热参数与名义动态控制参数的敏感性}
\label{F7-7}
\end{figure}

由图\ref{F7-7}(a)、(b)可见，内部导热能力降低时，端部与中部更难迅速交换热量，预冷温度场的相对非均匀程度增大。端面换热增强则同时提高外部散热强度并加快整体降温，因此固定预冷时刻的绝对温差不一定随换热系数单调增加；当电堆已经较接近环境温度时，可用于维持空间梯度的超温也随之减小。采用温差与热容加权均温超出环境温度之比，可更清楚地识别外部换热与内部导热竞争的影响。

启动敏感性表明，端板耦合导度对辅助电能的影响较突出。以完全冷却的名义方案为例，将端板耦合导度降低20\%后，辅助电能约降至2453.251 J；提高20\%后则增至约3190.957 J，显著偏离基准2855.408 J。端板耦合越强，启动过程中单电池向低温端板传递的热量越多，维持参考温升所需的辅助电能也越高。另一方面，个别控制参数扰动会导致未能完成测量保持停机，故图中的失败候选即使已消耗较少电能，也不能作为更优的成功方案。图\ref{F7-7}使用名义动态参数，其作用是揭示影响机制，不代表推荐方案已经在全部参数组合中得到严格鲁棒保证。

\begin{figure}[!h]
\centering
\includegraphics[width=.95\textwidth]{图13_名义与留裕度备选.pdf}
\caption{名义节能方案与推荐留裕度方案的能耗、时间及扰动表现}
\label{F7-8}
\end{figure}

由图\ref{F7-8}(a)、(b)可见，名义方案\(D_N\)三工况辅助电能为2855.408 J、0 J和1767.146 J，其中两个较冷工况的实际关热时刻均达到98.4 s，靠近允许的保持窗口上限。推荐方案\(D_R\)分别增加约74.588 J、293.609 J和88.239 J，将关热时刻提前至88.0 s、78.2 s和91.6 s。其目的在于为初温偏差、滤波滞后及物性失配保留可用余量，而不是在名义模型下进一步压低辅助能耗。

名义方案在原测量噪声与初场偏差试验中，仅20/90次满足全部可行性条件，成功样本集中于20 min工况；推荐方案在参数冻结后的独立噪声与初场试验中通过90/90次，并在24个物性失配与噪声联合场景中全部通过。两组噪声试验采用不同随机种子，因而这些比例可以说明各方案在所列场景下的表现，但不能作为严格配对的统计优势证明。固定恒功率\(C_H\)在相同最终验证场景下也全部通过，说明它本身具有较充分的启动时间余量。推荐动态方案的价值在于在这些已验证场景中保留可行性的同时降低辅助电能，而不是声称只有动态控制才能抵御扰动。

\begin{figure}[!h]
\centering
\includegraphics[width=.95\textwidth]{图16_观测器误差与测量停机.pdf}
\caption{参数失配与测量噪声下的观测误差及停机事件辨析}
\label{F7-9}
\end{figure}

图\ref{F7-9}进一步检验反馈信息的可获得性。图中采用未参与训练的参数失配与测量噪声样例，控制器仅使用独立观测模型和温度、电压采样。端板温度估计随物性失配出现偏差，说明观测状态并未直接复制被控对象的真实状态；冰风险估计与真实局部冰峰值之间也保留误差，不能视为对完整冰场的精确重构。以图\ref{F7-9}(c)的完全冷却样例为例，真实首次达标约为90.70 s，而滤波测温满足温度裕度并连续保持后，到95.0 s才确认关热；20 min样例的相应时刻约为82.61 s和91.0 s。这种差异来自传感采样、滤波和保持计时，说明实际关热时刻不能简单设为\(t_s+2\ \mathrm s\)。

因此，推荐方案的稳健性应由真实状态下的成功时刻、路径约束和关热前真实状态连续保持是否完成共同判定，不能仅以测量判据被触发为依据。图\ref{F7-8}(d)仍显示部分工况关热后回冷，也再次说明本节的扰动可行性验证与持续暖态验证属于不同层次。

\subsubsection{能量守恒、数值收敛与结果可靠性分析}

为解释节电的能量来源，将启动主窗口中的辅助加热、反应产热、相变净释热、环境散热和离散显热增量分别累计。按照与时间推进相同的热容和时间积分定义，离散能量残差写为
\begin{equation}
R_E=E_{\rm aux}+E_{\rm gen}+E_{\rm phase}-E_{\rm loss}-E_{\rm sensible}.
\tag{7-47}\label{eq:q4:energyresidual}
\end{equation}
其中\(E_{\rm phase}\)为带符号的相变净热，\(E_{\rm sensible}\)按各时间步使用的有效热容与温度增量累计。该量与离散方程保持一致，不将随含水状态变化的有效热容简单替换为一个常数乘总温升。

\begin{figure}[!h]
\centering
\includegraphics[width=.95\textwidth]{图09_能量收支与累积能耗.pdf}
\caption{推荐动态与固定恒功率策略的能量收支及累计辅助电能}
\label{F7-10}
\end{figure}

由图\ref{F7-10}(a)可见，各工况热输入总量与显热增量、环境散热之和相互吻合。完全冷却推荐方案在关热前输入辅助热2929.996 J、反应热1941.572 J和相变净热13.454 J，扣除环境散热272.961 J后，离散显热增量为4612.062 J。相同工况固定恒功率方案的反应热仅为194.959 J，而辅助输入为3260.281 J。因此，动态方案虽然延长了升温过程并增加环境散热，却通过更充分地利用反应热降低了所需辅助电能，同时使关热时端板温度较高。这也解释了其关热后回冷较固定策略有所减轻。

对于预冷20 min工况，推荐方案的反应热达到1533.862 J，明显高于293.609 J的辅助热；相变净热为\(-2.020\) J，环境散热为394.056 J，最终显热增量约为1431.395 J。预冷40 min时，辅助热和反应热分别为1855.385 J与2027.484 J，二者共同提供主要升温能量。图\ref{F7-10}(b)中，20 min工况累计辅助电能在末段转为水平，而其余两种工况仍随端部加热持续增加，与前述功率轨迹及提前停止实际加热的现象一致。

三工况推荐策略的离散能量残差绝对值约为\(10^{-11}\)--\(10^{-10}\) J，完整轨迹水量残差最大绝对值约为\(10^{-16}\)--\(10^{-15}\) kg，表明当前离散更新中的收支能够闭合。预冷阶段另行检查全堆内能减少与两端累计散热的一致性，并利用同网格矩阵指数解校验时间离散。守恒残差很小主要说明数值收支自洽，仍需结合独立参考和网格加密判断离散误差，不能单凭该残差推断物理预测没有误差。

\begin{figure}[!h]
\centering
\includegraphics[width=.95\textwidth]{图10_网格步长与控制周期检验.pdf}
\caption{预冷温度场与推荐动态启动结果的时间和空间收敛检验}
\label{F7-11}
\end{figure}

由图\ref{F7-11}(a)、(b)可见，预冷时间步从2 s减小至0.25 s时，20 min温度场相对同网格矩阵指数解的最大误差由约0.0166 K降至0.00208 K，继续减小至0.125 s后约为0.00104 K，表现出后向欧拉方法预期的收敛趋势。增加空间控制体数量后，七节点投影相对加密参考的误差也进一步下降，说明预冷初始场的离散精度能够满足三工况差异分析的需要。

对启动过程，固定正式控制参数和0.2 s采样周期，将每片水相控制体由58个加密至232个，同时把时间步从0.025 s减小至0.00625 s。三工况辅助电能分别由2929.996 J、293.609 J和1855.385 J变为2928.508 J、293.469 J和1854.436 J，相对变化均约为0.05\%；首次启动时间变化均小于0.008 s，实际关热时刻保持不变。图\ref{F7-11}(c)由此表明，主要能耗和时间结论对进一步联合加密较稳定。

图\ref{F7-11}(d)同时显示，局部冰峰值比总能耗更依赖空间分辨率。完全冷却工况由0.27725升至0.28658，20 min工况由0.08729升至0.08912，40 min工况由0.23957升至0.24623。尽管加密后的冰量仍远低于0.99约束，且不改变三工况的可行性判断，但其相对变化大于辅助电能的变化，因此不能将所有局部状态均描述为完全网格无关。综合基准复现、守恒检查、独立参考、离散加密与扰动验证，本章结果可支持给定模型及搜索范围内的性能比较；预冷和启动阶段沿用的有效物性、换热位置近似，以及尚缺新增实验校准的限制仍需保留。

\subsection{问题四小结}

针对不同预冷程度下电堆动态辅助启动问题，本文在前述水--冰--热--电耦合模型基础上，建立了整堆预冷传热模型，将预冷温度场映射为五片电池与两端端板的非均匀初态；进一步利用温度、电压及其变化率构造状态识别与独立观测器，通过前馈、PI反馈、温升不足补偿、安全增强和测量保持停机形成分片动态加热策略。以辅助电能为首要目标完成有限参数搜索后，再根据扰动训练场景选择留有温度和时间裕度的推荐方案。

结果表明，推荐方案在完全冷却、预冷20 min和预冷40 min工况下分别消耗2929.996 J、293.609 J和1855.385 J辅助电能，相对共同测量停机规则下的问题三固定功率方案节电10.13\%、68.26\%和17.28\%。其节能主要来自减少内部片持续加热，并利用较长启动过程中的反应热分担升温需求，代价是启动较慢且主窗口结冰水平较高；三工况电压、冰量及启动电荷仍满足约束。推荐方案通过所列独立扰动场景，但完全冷却与40 min工况关热后仍出现回冷，因此完成首次启动及短时保持并不等同于实现持续暖态。

在10--100 min预冷扫描中，19个离散工况均完成固定恒功率首次启动。随着预冷加深，初温降低，启动时间、辅助电能和过程最大片间温差增加，并趋近完全冷却结果；预冷场自身的绝对温差则逐渐减小。这说明初始储热量决定总体升温需求，而启动期间的片间一致性还受到端板热汇和分片热输入的影响。

% 第七章正文结束
\end{document}
